%% Consignes : « Toujours tres importante, la conclusion comporte : une
%% appreciation qualitative et quantitative des resultats et de l'avancement
%% du projet [...] des perspectives [...] un bilan personnel [...] Il est en
%% particulier demande a l'eleve de faire le lien avec le referentiel de
%% competences de l'ingenieur ESTIA. »

\subsection{Results Against the Objectives}
\label{sec:conclusion_objectifs}

Table~\ref{tab:bilan} assesses each objective of Section~\ref{sec:objectifs} against its
indicator. Four of five were met. The fifth was answered, and the answer was negative.

\begin{table}[htbp]
\centering
\caption{Objectives against outcomes}
\label{tab:bilan}
\small
\begin{tabularx}{\linewidth}{>{\raggedright\arraybackslash}p{0.7cm} >{\raggedright\arraybackslash}p{2.4cm} X}
\toprule
& \textbf{Status} & \textbf{Evidence} \\
\midrule
O1 & Met & Six-block aggregation with bidirectional propagation and an asymmetric adequacy
score, each term exactly decomposable; aggregation choice justified against the compensatory
alternative \\
O2 & Met & SupplyScore delivered: 73 modules, 1\,779 automated tests, multi-tenant
persistence, every write audited, operational and reloadable \\
O3 & Met & Propagation and criticality services operational, validated by dedicated
performance benchmarks at the scale of a real programme \\
O4 & Answered, negative & The signal does not predict recorded outcomes on this campaign. The
cause is identified: one block is structurally unable to respond to a shock \\
O5 & Met & Scientific report delivered; every reported figure regenerated from raw artefacts
by script rather than transcribed \\
\bottomrule
\end{tabularx}
\end{table}

\paragraph{Quantitatively.}
The delivered system covers the specified scope and is measurably robust: 1\,779 automated
tests, an auditable write path, and performance validated at the intended scale. The
experiment produced numbers rather than impressions: an area under the ROC curve between
$0.37$ and $0.56$, a Brier skill between $-3.4$ and $-6.9$, twenty-one confident forecasts of
which none materialised, and an eight-of-eight one-sided response at first exposure to the
displayed forecast.

\paragraph{Qualitatively.}
The mission set out to determine whether the declared-versus-real gap can be measured well
enough to act on. It cannot, yet, and the work establishes precisely why. A negative result
of this shape is more useful to the programme than a favourable score would have been,
because it is actionable: it names one component, states the mechanism, and specifies the
repair. A favourable score on a signal driven by a blind block would have been worse than
useless, since it would have been trusted.

\subsection{Limitations}
\label{sec:conclusion_limites}

Four limits bound what this report supports. The declarants were agents rather than human
consultants, so every statement about declaration behaviour applies to that population.
Results come from a single scenario, so the outcome base rate and event mix are properties of
one script. The positive class is very small, which bounds every discrimination claim in both
directions. And the carbon block was the least exercised of the six, for want of data to feed
it, which is the measurement gap noted in Section~\ref{sec:rse}.

\subsection{Perspectives}
\label{sec:perspectives}

\paragraph{For ALTEN.}
Four steps follow, in order. Repair the time block first, by making the calibrated events
write a lead-time impact alongside their existing effects; nothing else should be measured
until they do, since any campaign run against the current forecast layer will reproduce the
same result. Add a post-hoc calibration layer, fitted on data other than the evaluation
scenario. Make the forecast state its own ignorance, since a layer that cannot anticipate
exogenous shocks should not present a low probability in a confident numeric format. Then run
the campaign with human participants, which is the only run that can decide the hypotheses as
they were written, and which should happen only after the repair.

Beyond that, the adequacy layer is designed to be coupled to DISCO rather than to remain
standalone, and the network representation work carried out in parallel on the same programme
is the natural counterpart to connect it to.

\paragraph{For the author.}
The subject continues in the research thesis produced for Cranfield University, which treats
the same measurements at greater depth. The open question I would most like to pursue is the
one the treatment arm opened: what an advisory layer owes the person reading it when it does
not know.

\subsection{Personal Assessment}
\label{sec:bilan}

\subsubsection{What I learned}

The hardest problem of these six months was not mathematical. It was recognising, three
separate times and in three different disguises, that the system was measuring itself. Each
time the symptom looked like a modelling defect and each time the cause was circularity in
the measurement. I would now design the independence of two signals before designing either
of them.

The second lesson concerns what counts as a result. For a long stretch I treated a negative
outcome as a failure to be fixed before reporting it. The most useful findings in this report
are negative, and none would have been found by a project optimising for a positive one. The
related discipline, which the rejected dataset taught quickly, is that a confident
presentation is not evidence, and that the cheap verification step belongs before the
expensive one.

With hindsight I would change two things. I would have scored the forecast layer against
recorded outcomes far earlier, since doing so took a few hundred lines and revealed the
defect immediately where months of implementation had not. And I would have written the
outcome definition into the pre-registration alongside the thresholds, having since measured
how much it moves the answer.

\subsubsection{Link with the ESTIA engineering competency framework}

Table~\ref{tab:competences} maps the mission onto the competencies of the ESTIA framework
that it genuinely exercised. Appendix~\ref{app:competences} gives the detailed evidence.

\begin{table}[htbp]
\centering
\caption{Competencies exercised during the mission}
\label{tab:competences}
\small
\begin{tabularx}{\linewidth}{>{\raggedright\arraybackslash}p{1.1cm} >{\raggedright\arraybackslash}p{4.2cm} X}
\toprule
& \textbf{Competency} & \textbf{Evidence from this mission} \\
\midrule
CI1 & Analysis, synthesis, creativity, systemic and interdisciplinary approach & Combining
operations research, behavioural science, probabilistic modelling and software architecture
into one measurement \\
CI3 & Working in an international and intercultural context & Mission conducted and reported
in English, under a double-degree programme, alongside an English-speaking colleague on the
same research programme \\
CI5 & Identifying and solving complex or research problems using external resources & The
whole subject; and literally external, since the state of the art depended on university
library access unavailable within the company \\
CE3 & Leading and controlling a project against objectives, deadlines and risks & A four-touchpoint
weekly cadence, six phases, a documented mid-mission replanning, and a risk register of which
three items materialised \\
CE4 & Taking the socio-ecological transition into account & Carbon modelled as one of six
first-class blocks rather than as a correction, within a programme that is itself part of
ALTEN's environmental R\&D commitment \\
CE5 & Researching, selecting, qualifying and diffusing information & A candidate dataset
rejected on inspection before use; the same work written three ways for three different
audiences \\
CST1 & Defining an engineering strategy in a poorly known context & Converging a five-workstream
assignment onto one defensible contribution, and abandoning the first model when it proved
unable to represent the problem \\
CST2 & Modelling and simulating a system, capitalising the knowledge produced & The six-block
model, its propagation, the campaign simulation, and a reproducible derivation of every
reported figure \\
CST7 & Building models, methods and tools to aid engineering, in an R\&D perspective & The
delivered application and the experimental protocol, both designed to be continued by the
programme rather than to end with the placement \\
\bottomrule
\end{tabularx}
\end{table}

\subsubsection{What this mission clarified professionally}

I came into this placement expecting the interesting difficulty to be in the mathematics. It
was not. The mathematics behaved; what failed was the join between a model and the world it
was supposed to describe, and then the join between a model's output and the person meant to
act on it. Both failures were only visible because the work was measured against something
external rather than judged on its own internal consistency.

That is the kind of engineering I want to continue: the boundary where a model meets the
person who has to act on it. It is where the failures are least obvious and most consequential,
and this mission was a fairly direct demonstration of why.
